%! Author = lili
%! Date = 6/3/26


\section{Internal and general documentation}\label{sec:internal-and-general-documentation}
A lot of software solutions and methods for documentation exist.
I will not discuss the advantages and disadvantages of these, but instead mention general considerations.
The first and foremost being: Your documentation system has to be thorough, but easy to use.

The most sophisticated system is of no use, if half the team cannot use it on their device or it takes
weeks to learn it.

\subsection{Meeting protocols}\label{subsec:meeting-protocols}
In Bielefeld, we have a meeting template - that varies each year - for team meetings.
Team members are asked to prefill it.
This also opens up the option to put town tops anys time, lessening the amount of forgotten points.
It usually contains:
\begin{itemize}
    \item A ``Top'' section.
    Tops being discussion points where input or decisions are needed.
    \item An ``Info'' section that is not read out.
    Everyone is responsible for reading it themselves as it contains
    information relevant to the whole team.
    \item A ``Subteam update'' section that can later be used to lookup information.
\end{itemize}
Further information such as who was present or who wrote the protocol can be added.

\paragraph{Subteam meetings} are usually not as formally documented, but protocols are expected.

\paragraph{The minute taker} changes every meeting as we rotate through a list to render discussions about who has
to do it unnecessary.

\paragraph{It is advisable} to have meeting protocols from the beginning.
Depending on your team structure, experience and regarding time management, you should consider:
\begin{itemize}
    \item \textbf{Preparing agendas}: Starting meetings with a pre-prepared agenda outlining the topics for discussion.
    \item This helps keep the meeting focused and provides a structure for your notes.
    \item \textbf{Tracking decisions}: Document any decisions or agreements reached during the meeting to provide a clear record of team choices.
    \item \textbf{File Storage and Accessibility:} The meeting notes should be accessible to the whole team afterward and multiple versions should be most urgently avoided.
    \item \textbf{Summaries}: Some teams may find it helpful to create meeting summaries after the fact.
    Though this can also be a time drain.
\end{itemize}

\subsection{Task tracking}\label{subsec:task-tracking}
While there might not be as many tasks in the beginning I urge you to use some kind of tracking system.
Many apps are only available for certain browsers or OS, therefore we usually use an excel sheet.
I recommend writing down not only the task and the assignee, but also a deadline.

That way you can check in with each other, if tasks get lost.
Please remember to be kind with your team mates and that this is not about control or who is right, but about being
thorough.

\paragraph{Do this also for texts or other wiki content! }
In regard to texts and wiki content, we have an excel sheet containing:
\begin{itemize}
    \item Topic / Title
    \item Author(s)
    \item Wiki page it will be published on
    \item Category / Subteam
    \item Deadline
    \item Done?
    \item Three columns where other team members put their initials of they proofread the text
    \item A column where an advisor puts their initial if they proofread the text
    \item A checkbox for the author to check if they have reworked the text according to the proofreading and
    corrections / suggestions
    \item Is the text on the wiki?
\end{itemize}
This allows for automatic counting of how many texts / proofreading per person were done to avoid having like two
people do all the texts alone.

\subsection{Sponsoring}\label{subsec:sponsoring}
If you get sponsoring from companies, it is likely you promise them things in return or even make contracts.
For example presenting their logo on your wiki.

To keep contacts over the years, you need to not document the companies, but the contact data, who you spoke with
and possible additional info such as rules or recommendations they gave you.

To avoid looking unorganized and unprofessional, I recommend writing down who contacted the companies (if they have
been contacted at all) and the status of the outreach.

In Bielefeld, we use an excel sheet for all sponsoring contacts and a second sheet to document the won sponsors.

\noindent The general sheet contains:
\begin{itemize}
    \item Company name
    \item Person in charge
    \item Status
    \begin{itemize}
        \item First outreach (no answer yet)
        \item Second outreach (no answer yet)
        \item Not started
        \item Declined
        \item In contact
        \item Won
    \end{itemize}
    \item Contact person (if there is one)
    \item Phone number (if available)
    \item Email
    \item Additional info / notes (e.g.\ ``'', ``prefers high social media activity'')
    \item Products (they could possibly provide)
\end{itemize}
The sheet for won sponsors includes (besides company name, etc.):
\begin{itemize}
    \item Money sponsoring (if given)
    \item Commodity sponsoring (if given)
    \item Sponsoring status
    \begin{itemize}
        \item To be shipped
        \item Not transferred yet
        \item Received
    \end{itemize}
    \item Contract?
    \item Category (if you use this system)
    \begin{itemize}
        \item None
        \item Bronze
        \item Silver
        \item Gold
    \end{itemize}
    \item Did they provide the logo?
    \item A column for each benefit you can provide the companies where you note if they shall receive it and if you
    fulfilled that commitment.
    This can for example include:
    \begin{itemize}
        \item Logo on wiki
        \item Logo on team shirts
        \item Social media appreciation post
    \end{itemize}
\end{itemize}

\subsection{Images and videos}\label{subsec:images-and-videos}
Applications such as Microsoft Teams struggle with a high load of videos or images.
Consider using a cloud or shared
university file server.

I recommend to document if you have the consent of the people in the images and videos if you plan on releasing
footage of external people.

\paragraph{If your team has advisors or other people who stay for multiple years,} please consider having them meet
the sponsors.
This makes the contacts less dependent on the changing roster and makes it easier to secure funding independent of
the also changing projects topics.

\subsection{Attributions}\label{subsec:attributions}
Do not wait until the end to document your team attributions or external attributions.
Do so on a rolling basis and
make sure to explicitly distribute the responsibilities.
Otherwise, the subteam might end up thinking the
organization team does the attributions and the other way around, and they do not get done.

I recommend documenting them the way you have to submit them later on directly.

\subsection{HP contacts}\label{subsec:hp-contacts}
Similar to the sponsoring contact, document:
\begin{itemize}
    \item Name
    \item Institution / Affiliation
    \item Contact info
    \item Person in charge of contact
    \item Outreach status (``First outreach'', ``Declines'', etc.)
\end{itemize}
Additionally, document:
\begin{itemize}
    \item Stakeholder category
    \item Did they consent to you using their picture
    \item Did they consent to publication of interviews / transcripts or do they want to read ot first, ...
\end{itemize}

\paragraph{Consent Management} includes explaining how the ``freezing'' of iGEM Wikis works and that it will not be possible to change or remove information after the project ends.
Be mindful to receive feedback for quotes and transcripts\index{Transcripts} of interviews. \index{Interview}
Especially if you need to translate conversations to English, the stakeholders should get the chance to comment and, if necessary, correct your translations.

I am diving further into HP documentation in the section \nameref{sec:documentation-for-the-competition}.

\subsection{Staying on top of the documentation}\label{subsec:staying-on-top-of-the-documentation}
See also: \nameref{sec:internal-freeze-approach}
\subsection{Milestone approach}\label{subsec:milestone-approach}
In 2025, I tried the approach to set specific dates for check-ins with the team at the beginning of the wiki work,
akin to agile project management (but less micromanager-y, no stand-up meetings or daily scrums or whatever).
The team had specific design and documentation task that were to be finished and the expectation was that what was
not finished at the deadline had to be caught up on during co-working sessions.
After the design phase, biweekly checkins (and working sessions) were established with specifications in what kind
of product the various tasks should have.

We are now continuing this practice in 2026.

Tasks can be specific or abstract such as ``the first draft for every wiki text for HP work up
until this point is done'' or ``every text from the last intervall is proofread''.

Some abstract tasks we implemented are:
\begin{itemize}
    \item All existing (and proofread) texts (and graphics where applicable) from the last iteration are on the wiki (
    \textit{HTML / Template})
    \item Texts including sources exist for all past events and interviews (\textit{Word (texts) / BibTeX (sources)})
    \item All existing graphics and extra data have been uploaded to the upload tool (\textit{upload})
    \item Engineering cycle up to this point (outline) with Lab Team and including HP (\textit{Word or PowerPoint})
\end{itemize}
This works quite well and the team was usually up to date with their documentation.





